% =====================================================================
% 4. MODELING AND EVALUATION
% Rubric: models used, how the modelling was done, performance with
% appropriate metrics. Extra credit: discussion of choices and
% rationale, EXPLAINABLE AI, ITERATIVE IMPROVEMENT PROCESS.
% Source: experiments.csv / folds.csv (local), Roadmap Phases 0-4.
% Page budget: \todo{decide}
% =====================================================================
\section{Modeling and Evaluation}
\label{sec:modeling}

% --- 4.1 Evaluation protocol ----------------------------------------
\paragraph{Evaluation protocol.}
\todo{Repeated stratified k-fold (folds x repeats), metric =
\onenotch{} via a custom scorer in all CV and tuning, held-out test set
used once, mean $\pm$ std over folds.}

% --- 4.2 Ordinal target and decision rule ---------------------------
\paragraph{Using the ordinal structure.}
Because a one-class error is free, the class that maximises the
expected \onenotch{} is
$\hat k = \arg\max_k \big(p_{k-1} + p_k + p_{k+1}\big)$
rather than $\arg\max_k p_k$.
\todo{Compare: argmax vs window rule vs regression on the rank + rounding.}

% --- 4.3 Models compared --------------------------------------------
\begin{table}[h]
  \centering\small
  \caption{Model comparison under the same CV harness (mean $\pm$ std).}
  \label{tab:models}
  \begin{tabular}{lccc}
    \toprule
    Model & \onenotch{} & Exact acc. & MAE (rank) \\
    \midrule
    Majority class (baseline)      & \todo{} & \todo{} & \todo{} \\
    Logistic regression            & \todo{} & \todo{} & \todo{} \\
    k-nearest neighbours           & \todo{} & \todo{} & \todo{} \\
    Decision tree                  & \todo{} & \todo{} & \todo{} \\
    Random forest                  & \todo{} & \todo{} & \todo{} \\
    SVM                            & \todo{} & \todo{} & \todo{} \\
    Histogram gradient boosting    & \todo{} & \todo{} & \todo{} \\
    \bottomrule
  \end{tabular}
\end{table}

\paragraph{Tuning.}
\todo{Search method (e.g.\ halving random search), search space,
selected hyperparameters (full grid in the appendix).}

% --- 4.4 Iterative improvement (ablation) ---------------------------
% The rubric explicitly rewards this. Each row = one change, measured.
\paragraph{Iterative improvement.}
\begin{table}[h]
  \centering\small
  \caption{Ablation: cumulative effect of each step on CV \onenotch{}.}
  \label{tab:ablation}
  \begin{tabular}{clcc}
    \toprule
    \# & Step added & \onenotch{} & $\Delta$ \\
    \midrule
    0 & Baseline (probe model, basic preprocessing) & \todo{} & -- \\
    1 & \todo{e.g.\ window decision rule}           & \todo{} & \todo{} \\
    2 & \todo{e.g.\ target encoding high-card.}     & \todo{} & \todo{} \\
    3 & \todo{e.g.\ domain features}                & \todo{} & \todo{} \\
    4 & \todo{e.g.\ tuning}                         & \todo{} & \todo{} \\
    \midrule
    & Final model on held-out test set              & \todo{} & \\
    \bottomrule
  \end{tabular}
\end{table}

% --- 4.5 Error analysis ---------------------------------------------
\paragraph{Error analysis.}
\todo{Confusion matrix and share of errors by notch distance (0, 1, 2+);
which classes are confused.}

% --- 4.6 Explainability ---------------------------------------------
\paragraph{Explainability.}
\todo{Grouped permutation importance on held-out folds (per variable and
per feature group), partial dependence for the top drivers; link back to
objective O2.}

\begin{figure}[h]
  \centering
  \begin{subfigure}{0.49\linewidth}
    \centering\figplaceholder{Confusion matrix (ordered labels)}
    \caption{Confusion matrix.}
    \label{fig:confusion}
  \end{subfigure}\hfill
  \begin{subfigure}{0.49\linewidth}
    \centering\figplaceholder{Grouped permutation importance}
    \caption{Importance per feature group.}
    \label{fig:importance}
  \end{subfigure}
  \caption{Error analysis and explainability of the final model.}
\end{figure}
